\documentclass[11pt,a4paper]{article}

\usepackage[utf8]{inputenc}
\usepackage[T1]{fontenc}
\usepackage[margin=2.5cm]{geometry}
\usepackage{amsmath,amssymb,amsthm}
\usepackage{booktabs,array,longtable,tabularx,enumitem}
\usepackage{microtype}
\usepackage{xcolor}
\usepackage[colorlinks=true,
            linkcolor=blue!55!black,
            citecolor=blue!55!black,
            urlcolor=blue!55!black]{hyperref}
\usepackage{url}

\hypersetup{
  pdftitle={HypG v2.0.1: Normative Specification},
  pdfauthor={HypG Project},
  pdfsubject={Yellow Paper: algorithms, wire formats, and conformance contracts},
  pdfkeywords={post-quantum cryptography, Module-LWE, BLAKE3, ChaCha20, FORS, WOTS+, XMSS, specification}
}

\setlength{\emergencystretch}{2.5em}
\allowdisplaybreaks

\theoremstyle{plain}
\newtheorem{theorem}{Theorem}[section]
\newtheorem{proposition}[theorem]{Proposition}
\newtheorem{lemma}[theorem]{Lemma}
\newtheorem{corollary}[theorem]{Corollary}
\theoremstyle{definition}
\newtheorem{definition}[theorem]{Definition}
\newtheorem{assumption}[theorem]{Assumption}
\theoremstyle{remark}
\newtheorem{remark}[theorem]{Remark}

\newcommand{\concat}{\,\Vert\,}
\newcommand{\bits}{\{0,1\}}
\newcommand{\Z}{\mathbb{Z}}
\newcommand{\Ring}{\mathcal{R}_q}
\newcommand{\Hash}{\mathsf{H}}
\newcommand{\XOF}{\mathsf{XOF}}
\newcommand{\KDF}{\mathsf{KDF}}
\newcommand{\PRF}{\mathsf{PRF}}
\newcommand{\KeyGen}{\mathsf{KeyGen}}
\newcommand{\Encaps}{\mathsf{Encaps}}
\newcommand{\Decaps}{\mathsf{Decaps}}
\newcommand{\Encode}{\mathsf{Encode}}
\newcommand{\Decode}{\mathsf{Decode}}
\newcommand{\Pack}{\mathsf{Pack}}
\newcommand{\Unpack}{\mathsf{Unpack}}
\newcommand{\Sign}{\mathsf{Sign}}
\newcommand{\Verify}{\mathsf{Verify}}
\newcommand{\CBD}{\mathrm{CBD}}
\newcommand{\Compress}{\mathrm{Compress}}
\newcommand{\Decompress}{\mathrm{Decompress}}
\newcommand{\negl}{\mathsf{negl}}
\newcommand{\poly}{\mathsf{poly}}
\newcommand{\ppt}{\mathsf{PPT}}
\newcommand{\adv}{\mathcal{A}}
\newcommand{\ro}{\mathcal{O}}

\newenvironment{steps}[1]
  {\par\medskip\noindent\textbf{#1}\par\nopagebreak\small
   \begin{enumerate}[label=\arabic*:,leftmargin=*]}
  {\end{enumerate}\par\medskip}

\title{\textbf{HypG v2.0.1: Normative Specification}\\[0.4ex]
\large A Module-LWE KEM, Symmetric AEAD, and Stateless\\
Hash-Based Signature Scheme\\[0.3ex]
\normalsize Yellow Paper --- Algorithms, Wire Formats, and Conformance Contracts\\[0.2ex]
\small Workspace 2.0.0; KEM/Symm/Hash 2.0.0; DS 2.0.1}

\author{HypG Project\\
\small Independent research prototype}
\date{October 2026}

\begin{document}
\maketitle

\begin{abstract}
This document is the normative specification of the HypG research
prototype.  It defines the byte-level encodings, algorithms, domain
separators, API contracts, and frozen test-vector inventory for the
Module-LWE KEM, the BLAKE3 hash/XOF/keyed-hash layer, the
ChaCha20 plus keyed-BLAKE3 authenticated-encryption scheme (one-shot
and streaming), and the stateless hash-based signature scheme.  The
workspace version is 2.0.0; the signature wire/spec revision is
separately identified as \textbf{2.0.1} because of two changes: the
explicit one-byte $\PRF$/$F$ domain tags and FIPS-205-style promotion of
odd WOTS+ L-tree nodes.

The document is normative for behaviour, not for security.  It does not
assert IND-CCA2, UF-CMA, SUF-CMA, or AEAD security for the exact
constructions.  It does not claim NIST Category-5 security, FIPS
conformance, certification, constant-time machine code, or side-channel
resistance.  In case of disagreement, the Rust implementation and the
frozen KAT files take precedence over this prose.
\end{abstract}

\noindent\textbf{Keywords:} specification; Module-LWE; Fujisaki--Okamoto;
implicit rejection; BLAKE3; ChaCha20; Encrypt-then-MAC; FORS; WOTS+;
XMSS; hypertree; safe Rust; known-answer tests.

\medskip
\noindent\textbf{Document status.} Normative technical report for the
v2.0 artifact.  It is not peer-reviewed, not an audit, and not a
security certification.

\tableofcontents
\newpage

%======================================================================
\section{Scope, status, and normative precedence}
\label{sec:scope}

\subsection{Document identity}

\begin{table}[t]
\centering
\small
\begin{tabular}{@{}llll@{}}
\toprule
Artifact & Code version & Wire/spec revision & Frozen vectors \\
\midrule
Workspace & 2.0.0 & --- & --- \\
KEM & 2.0.0 & KEM v2.0 & \texttt{kats/v2.0.0.txt} \\
Symm / AEAD & 2.0.0 & Symm v2.0 & \texttt{kats-symm/v2.0.0.txt} \\
Hash API & 2.0.0 & Hash v2.0 & \texttt{kats-hash/v2.0.0.txt} \\
Signatures & 2.0.0 & \textbf{DS v2.0.1} & \texttt{kats-ds/v2.0.1.txt} \\
\bottomrule
\end{tabular}
\caption{Version identity.  The literal \texttt{HYP-G/v07/...} strings
are frozen domain separators, not a statement that the current
parameter set is v0.7.}
\label{tab:versions}
\end{table}

\subsection{Normative language}

The key words MUST, MUST NOT, REQUIRED, SHALL, SHOULD, SHOULD NOT, MAY,
and OPTIONAL are to be interpreted as described in RFC~2119~\cite{rfc2119}
when they appear in this document.  In particular:

\begin{itemize}[leftmargin=*]
\item A conforming implementation MUST produce the bytes specified for
  every algorithm in this document.
\item A verifier MUST reject a malformed encoding as specified; it MUST
  NOT silently accept an alternative encoding.
\item A caller MUST satisfy the caller-side preconditions (fresh seeds,
  fresh encapsulation messages, unique nonces, bounded signing volume,
  and secret storage).
\end{itemize}

\subsection{Precedence}

The final arbiters of behaviour are:

\begin{enumerate}[leftmargin=*]
\item the Rust code in \texttt{crates/};
\item the active frozen KAT files in \texttt{crates/hypg-kat/};
\item the normative algorithms and tables in this document.
\end{enumerate}

If the prose and the code disagree, the code and KATs win; the prose is
then defective and must be corrected.  This is intentional: the
specification is a research artifact, and silent divergence between
paper and implementation is a bug.

\subsection{Out of scope}

This document does not specify a complete protocol.  It does not define
key storage, RNG construction, nonce assignment, message framing, replay
protection, identity binding, key confirmation, key rotation, or
authorization.  Those are caller responsibilities.  The document also
does not provide a security proof; see Section~\ref{sec:security}.

%======================================================================
\section{Notation and encodings}
\label{sec:notation}

\subsection{Bytes and integers}

\begin{definition}[Byte string]
A byte is an element of $\{0,\ldots,255\}$.  A byte string is a finite
sequence of bytes; $\bits^n$ denotes bit strings of length $n$, and
$\bits^*$ finite bit strings.
\end{definition}

\begin{definition}[Little-endian encoding]
For an integer $x$ with $0\le x<2^t$, $\mathrm{LE}_t(x)$ is the
$t/8$-byte encoding of $x$ with the least-significant byte first.  For
example, $\mathrm{LE}_{32}(1)=01\,00\,00\,00$.
\end{definition}

Unless stated otherwise, every multi-byte integer in this specification
is little-endian.  Byte ranges are half-open: $x[a:b]$ denotes bytes
$a$ through $b-1$.  The symbol $\concat$ denotes concatenation.  No
implicit padding is ever inserted.

\subsection{Rings and vectors}

\begin{definition}[Ring]
Let
\[
  \Ring=\Z_q[X]/(X^{256}+1), \qquad q=4096=2^{12}.
\]
An element $a\in\Ring$ is represented by its coefficient vector
$(a_0,\ldots,a_{255})$ with $a_i\in[0,q)$.  Addition and subtraction
are coefficient-wise modulo $q$.
\end{definition}

\begin{definition}[Negacyclic multiplication]
For $a,b\in\Ring$,
\[
 (a*b)_k=\sum_{\substack{i+j=k\\0\le i,j<256}}a_i b_j
        -\sum_{\substack{i+j=k+256\\0\le i,j<256}}a_i b_j
 \pmod q.
\]
\end{definition}

The implementation evaluates the two sums directly with an integer
accumulator and reduces modulo $q$ with a 12-bit mask.  The loop bounds
are public.  For module rank $k=4$, vectors are written
$a=(a_0,\ldots,a_{k-1})$ and
$\langle a,b\rangle=\sum_{i=0}^{k-1}a_i b_i$.

\subsection{The CBD distribution}

\begin{definition}[$\CBD(2)$]
On 128 uniformly random bytes, let $b_0,b_1,b_2,b_3$ be the four
successive bits for a coefficient.  Define
\[
 v=(b_0+b_1)-(b_2+b_3).
\]
The coefficient is $v$ mapped to $[0,q)$ by adding $q$ and masking.
\end{definition}

\begin{lemma}[CBD moments]
\label{lem:cbd}
For $\CBD(2)$, the support is $\{-2,-1,0,1,2\}$ and
\[
 \Pr[-2]=1/16,\quad \Pr[-1]=4/16,\quad \Pr[0]=6/16,\quad
 \Pr[1]=4/16,\quad \Pr[2]=1/16.
\]
The mean is $0$ and the variance is $1$.
\end{lemma}

\begin{proof}
The distribution is the difference of two independent
$\mathrm{Binomial}(2,1/2)$ variables.  Its probability generating
function is $(1/4+1/2z+1/4z^2)(1/4+1/2z^{-1}+1/4z^{-2})$, whose
coefficients are the probabilities above.  The mean is the difference of
the two binomial means, and the variance is the sum of the two
variances $2\cdot(1/2)(1/2)=1/2$.
\end{proof}

%======================================================================
\section{Core primitive: CA-PRF}
\label{sec:caprf}

\subsection{State and round function}

\begin{definition}[CA state and rule]
A CA state is a vector $S\in(\Z_{2^{64}})^{32}$, serialized as 32
little-endian 64-bit words (256 bytes).  A CA rule is a vector
$R\in(\Z_{2^{64}})^8$, serialized as eight little-endian 64-bit words
(64 bytes).  The public constants are
\[
 \begin{aligned}
 \mathrm{RC\_BASE}&=0x9E3779B97F4A7C15,\\
 \mathrm{RC\_STRIDE}&=0xBF58476D1CE4E5B9,\\
 \mathrm{DIFF\_MUL}&=0xD6E8FEB86659FD93.
 \end{aligned}
\]
\end{definition}

\begin{definition}[One CA round]
For round index $t\in[0,32)$ and word index $i\in[0,32)$, compute, with
all arithmetic modulo $2^{64}$ and $\mathrm{ROTL}$ a 64-bit left
rotation:
\[
\begin{aligned}
 rc &= \mathrm{RC\_BASE}+(t+1)\cdot\mathrm{RC\_STRIDE},\\
 a &= S_i,\quad b=S_{(i+1)\bmod 32},\quad
 c=S_{(i+5)\bmod 32},\quad d=S_{(i+17)\bmod 32},\\
 rk &= R_{i\bmod 8}\oplus rc\oplus
       (i\cdot\mathrm{RC\_STRIDE}),\\
 t_1 &= a\oplus\mathrm{ROTL}(b,13)\oplus\mathrm{ROTL}(c,7),\\
 t_2 &= (t_1\land d)\oplus((\lnot t_1)\land b),\\
 t_3 &= t_2+(rk\oplus(i+t)),\\
 t_4 &= t_3\oplus\mathrm{ROTL}(t_3,29),\\
 t_5 &= (t_4\cdot\mathrm{DIFF\_MUL})\oplus\mathrm{ROTL}(d,11),\\
 S'_i &= t_5\oplus\mathrm{ROTL}(S_{(i+16)\bmod 32},17)\oplus rc.
\end{aligned}
\]
The round is applied 32 times.
\end{definition}

All loops have public bounds.  The temporary array is zeroized after
each round.  The implementation contains no secret-dependent branch or
index in the round.

\subsection{CA-PRF definition}

\begin{definition}[CA-PRF]
For a 32-byte seed $x$, public domain $d$, and output length $L$:
\[
\begin{aligned}
 \mathit{init} &= \XOF(x\concat d\concat\texttt{"-init"},256),\\
 \mathit{rule} &= \XOF(x\concat d,64),\\
 S_0 &= \mathrm{LE}^{-1}_{64}(\mathit{init}),\\
 S_{32} &= \mathrm{Round}^{32}(S_0;\mathit{rule}),\\
 \mathit{out} &= \XOF(d\concat\mathrm{LE}_{64}(S_{32})
   \concat d\concat\texttt{"-out"},L).
\end{aligned}
\]
\end{definition}

\begin{remark}[Domain truncation]
The implementation copies at most 59 bytes of $d$ before
\texttt{"-init"} and at most 60 bytes before \texttt{"-out"}.  Every
protocol domain in Section~\ref{sec:domains} is shorter than these
limits, so no truncation occurs.  The rule expansion uses the full $d$.
\end{remark}

\begin{remark}[Security role of the CA]
The final BLAKE3 XOF is the security-relevant expander.  If BLAKE3 is a
secure XOF or random oracle, the output is pseudorandom regardless of the
CA.  The CA is a deterministic mixing and domain-separation stage, not
an independently proven hardness assumption.
\end{remark}

\begin{lemma}[Prefix consistency]
\label{lem:caprf-prefix}
For fixed $(x,d)$ and $L_1\le L_2$, the first $L_1$ bytes of
$\mathrm{CA\text{-}PRF}(x,d,L_2)$ equal
$\mathrm{CA\text{-}PRF}(x,d,L_1)$.
\end{lemma}

\begin{proof}
The state and all preceding stages are independent of $L$.  The final
XOF is prefix-consistent by definition.
\end{proof}

%======================================================================
\section{Hash, XOF, and KDF layer}
\label{sec:hash}

\subsection{Functions}

\begin{definition}[Hash API]
Let $\Hash(x)$ be unkeyed BLAKE3-256, $\Hash_k(x)$ keyed BLAKE3-256
with 32-byte key $k$, and $\XOF(x,L)$ the first $L$ bytes of the
BLAKE3 XOF.  Define
\[
\begin{aligned}
 \mathrm{expand}(x,d,L) &= \XOF(x\concat d,L),\\
 \mathrm{kdf}(d,\mathit{parts},L) &=
   \XOF(d\concat\mathit{parts}_0\concat\mathit{parts}_1\concat\cdots,L),\\
 \mathrm{derive}(c,k) &= \mathrm{BLAKE3.derive\_key}(c,k).
\end{aligned}
\]
\end{definition}

To avoid a proliferation of symbols, $\Hash$ is overloaded: it denotes
the unkeyed one-argument hash in the KEM/hash layer and the keyed
four-argument tweakable hash in the DS layer; the argument count and
context disambiguate.

BLAKE3 is the hash primitive of the suite~\cite{blake3}.  The KEM uses
$\mathrm{expand}$ only through CA-PRF.  The KEM shared-secret derivation
uses $\mathrm{kdf}$ with the domain
\texttt{HYP-G/v07/ss} and the fixed-length parts $(m,ct,\Hash(pk))$.
The rejection derivation uses the fixed-length parts
$(\mathit{fake},ct)$.

\begin{lemma}[XOF prefix property]
\label{lem:xof-prefix}
$\Hash(x)$ equals the first 32 bytes of $\XOF(x,64)$.
\end{lemma}

\begin{proof}
This is the BLAKE3 root-bytes property and is pinned by the repository
tests for empty, short, chunk-edge, and multi-chunk inputs.
\end{proof}

\subsection{Framing requirements}

Neither $\mathrm{expand}$ nor $\mathrm{kdf}$ inserts separators or
lengths.  A caller MUST ensure that its part sequence is injective.
This specification uses fixed-length parts for every KDF call.  A new
caller with variable-length parts MUST add explicit lengths or use
\texttt{derive\_key}.

\subsection{Streaming hasher}

\texttt{Hasher} supports \texttt{new}, \texttt{new\_keyed},
\texttt{update}, \texttt{finalize32}, and \texttt{finalize\_xof}.  The
first 32 bytes of \texttt{finalize\_xof} equal \texttt{finalize32}.
The implementation uses a fixed 54-level chaining-value stack and does
not allocate.

%======================================================================
\section{Polynomial encodings and compression}
\label{sec:poly}

\subsection{Message encoding}

\begin{definition}[Encode/Decode]
For a 32-byte message $m$, let
\[
 \Encode(m)_i=2048\cdot((m_{\lfloor i/8\rfloor}\gg(i\bmod 8))\land1).
\]
For $v\in\Ring$, let
\[
 \Decode(v)_i=
 \left(\left(\left((v_i+1024)\bmod q\right)\gg11\right)\land1\right),
\]
with the bits packed least-significant first.
\end{definition}

\subsection{Compression}

\begin{definition}[10-bit and 5-bit compression]
For a coefficient $c\in[0,q)$:
\[
 \begin{aligned}
 \Compress_{10}(c)&=((c+2)\gg2)\land0x3FF,\\
 \Decompress_{10}(c)&=(c\land0x3FF)\ll2,\\
 \Compress_5(c)&=((c+64)\gg7)\land0x1F,\\
 \Decompress_5(c)&=(c\land0x1F)\ll7.
 \end{aligned}
\]
\end{definition}

\begin{lemma}[Residual ranges]
\label{lem:residual}
For $c\in[0,q)$, write
$\Delta_1=\Decompress_{10}(\Compress_{10}(c))-c\pmod q$ and
$\Delta_2=\Decompress_5(\Compress_5(c))-c\pmod q$.  Interpreted in the
signed range,
\[
 \Delta_1\in\{-1,0,1,2\},\qquad
 \Delta_2\in\{-63,-62,\ldots,64\}.
\]
\end{lemma}

\begin{proof}
Write $c=4a+r$ with $0\le r<4$.  Then
$\Compress_{10}(c)=a$ for $r\le1$ and $a+1$ for $r\ge2$, giving
$\Delta_1=-r$ or $4-r$.  For the 5-bit case, write $c=128a+r$ with
$0\le r<128$; the quotient is $a$ for $r<64$ and $a+1$ for $r\ge64$,
giving $\Delta_2=-r$ or $128-r$.
\end{proof}

\subsection{Packing}

\begin{definition}[12-bit packing]
Two coefficients $c_0,c_1\in[0,4095]$ are packed into three bytes:
\[
 b_0=c_0\land0xFF,\quad
 b_1=((c_0\gg8)\land0x0F)\lor((c_1\land0x0F)\ll4),\quad
 b_2=(c_1\gg4)\land0xFF.
\]
\end{definition}

The 10-bit and 5-bit packers use a little-endian bitstream:
coefficient $i$ occupies bits $[i\cdot w,(i+1)\cdot w)$ for
$w\in\{10,5\}$.

\begin{lemma}[Packing bijectivity]
\label{lem:pack}
$\Pack_{12}/\Unpack_{12}$, $\Pack_{10}/\Unpack_{10}$, and
$\Pack_5/\Unpack_5$ are inverse bijections on their fixed-size arrays.
\end{lemma}

\begin{proof}
The 12-bit packing is a bijection because the three bytes contain
exactly the 24 bits of the two 12-bit values in a fixed order.  The
10-bit and 5-bit packers write each coefficient into a fixed bit range
of a little-endian bitstream of exactly $256w$ bits; no overlapping
writes and no unused prefix bits occur.
\end{proof}

%======================================================================
\section{KEM specification}
\label{sec:kem}

\subsection{Parameters and wire sizes}

\begin{table}[t]
\centering
\small
\begin{tabular}{@{}lrl@{}}
\toprule
Symbol & Value & Meaning \\
\midrule
$n$ & 256 & ring degree \\
$k$ & 4 & module rank; dimension 1024 \\
$q$ & $2^{12}=4096$ & modulus \\
$\eta$ & 2 & $\CBD(2)$ for $s,e,r,e_1,e_2$ \\
$m$ & 32 bytes & FO message; one bit per coefficient \\
$pk$ & 1568 bytes & $\mathit{seed}_A(32)\concat\Pack_{12}(b)(1536)$ \\
$ct$ & 1440 bytes & $\Pack_{10}(c_1)(1280)\concat\Pack_5(c_2)(160)$ \\
$sk$ & 1632 bytes & $\mathit{seed}(32)\concat\mathit{fake}(32)\concat pk$ \\
$ss$ & 32 bytes & XOF output \\
\bottomrule
\end{tabular}
\caption{KEM parameters and sizes.}
\label{tab:kemparams}
\end{table}

\subsection{Matrix and vector expansion}

\begin{definition}[ParseMatrix]
The CA-PRF output of length $4\cdot4\cdot384=6144$ bytes is read in
row-major order.  Entry $(i,j)$ is the 384-byte block at offset
$(4i+j)\cdot384$, interpreted as 256 12-bit coefficients by
$\Unpack_{12}$.
\end{definition}

\begin{definition}[ParseCBDVector]
A 512-byte CA-PRF output is split into four 128-byte blocks; each block
is mapped through $\CBD(2)$ to one polynomial.
\end{definition}

\subsection{Key generation}

\begin{steps}{Algorithm 1: $\KeyGen(\mathit{seed})$}
\item Input: a 32-byte seed with at least 256 bits of entropy.
\item $\mathit{seed}_A\gets\mathrm{CA\text{-}PRF}(\mathit{seed},
  \texttt{"HYP-G/v07/seed-a"},32)$.
\item $A\gets\mathrm{ParseMatrix}(
  \mathrm{CA\text{-}PRF}(\mathit{seed}_A,
  \texttt{"HYP-G/v07/poly-a"},6144))$.
\item $s\gets\mathrm{ParseCBDVector}(
  \mathrm{CA\text{-}PRF}(\mathit{seed},
  \texttt{"HYP-G/v07/sec-s"},512))$.
\item $e\gets\mathrm{ParseCBDVector}(
  \mathrm{CA\text{-}PRF}(\mathit{seed},
  \texttt{"HYP-G/v07/sec-e"},512))$.
\item For $i=0,\ldots,3$: $b_i\gets
  \sum_{j=0}^{3}A_{i,j}*s_j+e_i$ in $\Ring$.
\item $pk\gets\mathit{seed}_A\concat\Pack_{12}(b_0)\concat\cdots
  \concat\Pack_{12}(b_3)$.
\item $\mathit{fake}\gets\mathrm{CA\text{-}PRF}(\mathit{seed},
  \texttt{"HYP-G/v07/fake"},32)$.
\item $sk\gets\mathit{seed}\concat\mathit{fake}\concat pk$.
\item Return $(pk,sk)$.  Best-effort erase intermediate secret buffers.
\end{steps}

\subsection{Encapsulation}

\begin{steps}{Algorithm 2: $\Encaps(pk,m)$}
\item Input: a 1568-byte public key and a 32-byte FO message $m$.
  The caller MUST supply a fresh CSPRNG output for $m$ on every call.
\item Parse $pk$ as $(\mathit{seed}_A,b_0,\ldots,b_3)$.
\item $A\gets\mathrm{ParseMatrix}(
  \mathrm{CA\text{-}PRF}(\mathit{seed}_A,
  \texttt{"HYP-G/v07/poly-a"},6144))$.
\item $\mathit{rseed}\gets
  \Hash(m\concat pk\concat\texttt{"HYP-G/v07/r-seed"})$.
\item $r\gets\mathrm{ParseCBDVector}(
  \mathrm{CA\text{-}PRF}(\mathit{rseed},
  \texttt{"HYP-G/v07/rand-r"},512))$.
\item $e_1\gets\mathrm{ParseCBDVector}(
  \mathrm{CA\text{-}PRF}(\mathit{rseed},
  \texttt{"HYP-G/v07/rand-e1"},512))$.
\item $e_2\gets\CBD(\mathrm{CA\text{-}PRF}(\mathit{rseed},
  \texttt{"HYP-G/v07/rand-e2"},128))$.
\item For $i=0,\ldots,3$: $c_{1,i}\gets
  \sum_{j=0}^{3}A_{j,i}*r_j+e_{1,i}$.
\item $c_2\gets\sum_{i=0}^{3}b_i*r_i+e_2+\Encode(m)$.
\item $ct\gets\Pack_{10}(\Compress_{10}(c_{1,0}))\concat\cdots
  \concat\Pack_{10}(\Compress_{10}(c_{1,3}))
  \concat\Pack_5(\Compress_5(c_2))$.
\item $ss\gets\XOF(\texttt{"HYP-G/v07/ss"}\concat m\concat ct
  \concat\Hash(pk),32)$.
\item Return $(ct,ss)$.
\end{steps}

\subsection{Decapsulation}

\begin{steps}{Algorithm 3: $\Decaps(sk,ct)$}
\item Parse $sk$ as $(\mathit{seed},\mathit{fake},pk)$.
\item Parse $pk$ and expand $A$ as in encapsulation.
\item $s\gets\mathrm{ParseCBDVector}(
  \mathrm{CA\text{-}PRF}(\mathit{seed},
  \texttt{"HYP-G/v07/sec-s"},512))$.
\item Unpack and decompress $ct$ into $(c_1,c_2)$.
\item $v\gets c_2-\sum_{i=0}^{3}s_i*c_{1,i}$.
\item $m'\gets\Decode(v)$.
\item Recompute $ct'$ exactly as in Algorithm~2 from $(pk,m')$.
\item $\mathit{ok}\gets\mathrm{ct\_eq}(ct',ct)$ (full scan, returns
  $1$ on equality and $0$ otherwise).
\item $\mathit{pk\_hash}\gets\Hash(pk)$.
\item $ss_{\mathrm{ok}}\gets\XOF(
  \texttt{"HYP-G/v07/ss"}\concat m'\concat ct
  \concat\mathit{pk\_hash},32)$.
\item $ss_{\mathrm{rej}}\gets\XOF(
  \texttt{"HYP-G/v07/reject"}\concat\mathit{fake}\concat ct,32)$.
\item $ss\gets\mathrm{ct\_select}(\mathit{ok},ss_{\mathrm{ok}},
  ss_{\mathrm{rej}})$.
\item Return $ss$.  The function never returns an error.
\end{steps}

\subsection{Wire layout}

\begin{table}[t]
\centering
\small
\begin{tabular}{@{}rrl@{}}
\toprule
Offset & Length & Field in $pk$ \\
\midrule
0 & 32 & $\mathit{seed}_A$ \\
32 & 384 & packed $b_0$ \\
416 & 384 & packed $b_1$ \\
800 & 384 & packed $b_2$ \\
1184 & 384 & packed $b_3$ \\
\bottomrule
\end{tabular}
\hfill
\begin{tabular}{@{}rrl@{}}
\toprule
Offset & Length & Field in $ct$ \\
\midrule
0 & 320 & packed compressed $c_{1,0}$ \\
320 & 320 & packed compressed $c_{1,1}$ \\
640 & 320 & packed compressed $c_{1,2}$ \\
960 & 320 & packed compressed $c_{1,3}$ \\
1280 & 160 & packed compressed $c_2$ \\
\bottomrule
\end{tabular}
\caption{KEM public-key and ciphertext layouts.  The secret key is
$\mathit{seed}(32)\concat\mathit{fake}(32)\concat pk(1568)$.}
\label{tab:kemwire}
\end{table}

\subsection{Correctness lemmas}

\begin{lemma}[KEM noise identity]
\label{lem:kem-noise}
For a ciphertext produced by Algorithm~2 and unpacked by Algorithm~3,
before compression effects,
\[
 c_2-\langle s,c_1\rangle
 =\Encode(m)+\langle e,r\rangle+e_2-\langle s,e_1\rangle.
\]
\end{lemma}

\begin{proof}
Expand $b=A s+e$ and $c_1=A^T r+e_1$, then use associativity and
commutativity of multiplication in $\Ring$.  The mixed terms
$\langle A s,r\rangle$ and $\langle s,A^T r\rangle$ are the same sum
over the same index pairs.
\end{proof}

\begin{lemma}[Compression residual identity]
\label{lem:kem-residual}
Let $\Delta_{1,i}$ be the residual on $c_{1,i}$ and $\Delta_2$ the
residual on $c_2$.  Then
\[
 v'=\Encode(m)+\langle e,r\rangle+e_2-\langle s,e_1\rangle
 +\Delta_2-\langle s,\Delta_1\rangle.
\]
\end{lemma}

\begin{proof}
Substitute $c_{1,i}'=c_{1,i}+\Delta_{1,i}$ and
$c_2'=c_2+\Delta_2$ into $v'=c_2'-\langle s,c_1'\rangle$ and apply
Lemma~\ref{lem:kem-noise}.
\end{proof}

\begin{lemma}[Deterministic re-encryption]
\label{lem:kem-reenc}
If $\Decode(v')=m$, then Algorithm~3 recomputes $ct'=ct$ and selects
$ss_{\mathrm{ok}}$.
\end{lemma}

\begin{proof}
The expansions, PKE encryption, compression, and packing are
deterministic functions of their inputs.  With $m'=m$, every
intermediate value equals the encapsulation value, including the
fixed-length packed ciphertext.  The equality test therefore returns
$1$.
\end{proof}

\paragraph{Failure probability.}
The working variance model is
\[
 \mathrm{Var}[\langle e,r\rangle]=1024,\quad
 \mathrm{Var}[\langle s,e_1\rangle]=1024,\quad
 \mathrm{Var}[e_2]=1,
\]
\[
 \mathrm{Var}[\langle s,\Delta_1\rangle]=1280,\quad
 \mathrm{Var}[\Delta_2]=1365.25,
\]
for a total of $4694.25$ and $\sigma\approx68.5$.  Rounding up to
$\sigma\approx70$ gives a working margin of approximately
$14.6\sigma$ against the 1024 decision boundary.  This is a heuristic,
not a formal tail bound.  No analytic $2^{-128}$ failure bound is
claimed, and the 32-trial unit test is a regression tripwire.

\subsection{Caller semantics}

The 32-byte \texttt{coins} argument of the API is $m$, not a
domain-separated auxiliary nonce.  It is embedded in $c_2$ and is the
input to the FO randomness generator.  It MUST be fresh and uniformly
random for every encapsulation under a given public key.  The KEM has
no internal RNG; key generation, encapsulation, and decapsulation are
pure deterministic functions of their inputs.

\subsection{Malformed inputs}

The public-key and ciphertext parsers are infallible fixed-size
wrappers.  They do not validate that a public key was produced by
Algorithm~1, and they do not reject non-canonical algebraic structure
beyond unpacking 12-bit fields.  A malformed but correctly sized input
cannot cause an out-of-bounds access in safe Rust; decapsulation is
total.  It can of course produce a useless shared secret.

%======================================================================
\section{AEAD specification}
\label{sec:aead}

\subsection{Parameters}

\begin{table}[t]
\centering
\small
\begin{tabular}{@{}lr@{}}
\toprule
Parameter & Value \\
\midrule
Master key & 32 bytes \\
Encryption subkey & 32 bytes \\
MAC subkey & 32 bytes \\
Nonce & 12 bytes \\
Tag & 32 bytes \\
ChaCha20 block & 64 bytes, 20 rounds \\
Counter & 32 bits, initial 0 \\
Maximum message & $2^{38}$ bytes ($2^{32}$ blocks) \\
\bottomrule
\end{tabular}
\caption{AEAD parameters.}
\label{tab:aeadparams}
\end{table}

\subsection{Key schedule}

\[
 k_e=\mathrm{derive}(\texttt{"HYP-G/v07/symm-enc"},K),\qquad
 k_m=\mathrm{derive}(\texttt{"HYP-G/v07/symm-mac"},K).
\]
ChaCha20 is the IETF 20-round ARX stream cipher of RFC~8439~\cite{rfc8439}.  The master key is not used directly for ChaCha20 or for the MAC.

\subsection{One-shot sealing and opening}

\begin{steps}{Algorithm 4: $\mathrm{Seal}(K,N,AD,M)$}
\item Require $|C|=|M|$ and $|M|\le2^{38}$; otherwise erase $C$ and
  $T$ and return $\bot$.
\item $(k_e,k_m)\gets$ key schedule.
\item $C\gets M$; apply ChaCha20 with key $k_e$, nonce $N$, and
  counter $0$.
\item $T\gets\mathrm{KeyedBLAKE3}(k_m,
  N\concat AD\concat C\concat\mathrm{LE}_{64}(|AD|)
  \concat\mathrm{LE}_{64}(|C|))$.
\item Return $(C,T)$.
\end{steps}

\begin{steps}{Algorithm 5: $\mathrm{Open}(K,N,AD,C,T)$}
\item Require $|P|=|C|$ and $|C|\le2^{38}$; otherwise return $\bot$.
\item $(k_e,k_m)\gets$ key schedule.
\item $T'\gets\mathrm{KeyedBLAKE3}(k_m,
  N\concat AD\concat C\concat\mathrm{LE}_{64}(|AD|)
  \concat\mathrm{LE}_{64}(|C|))$.
\item If $\mathrm{ct\_eq}(T',T)\ne1$, return $\bot$ without writing
  $P$.
\item $P\gets C$; apply ChaCha20 with key $k_e$, nonce $N$, and
  counter $0$.
\item Return $P$.
\end{steps}

The length suffix is mandatory even when $AD$ or $C$ is empty.  The
tag covers the ciphertext, not the plaintext.

\subsection{Streaming state machine}

\texttt{StreamSealer} and \texttt{StreamOpener} maintain: a 64-byte
keystream block, a partial-block offset, a 64-bit next-counter value, a
running keyed-BLAKE3 state, the AD length, the ciphertext length, and a
failure flag.  \texttt{new} absorbs $N\concat AD$.  Sealer
\texttt{update} copies plaintext, XORs the next keystream bytes, adds
the chunk length using checked addition, and absorbs the ciphertext.
Opener \texttt{update} copies ciphertext, XORs tentatively, adds the
length, and absorbs the ciphertext.

\begin{proposition}[Streaming identity]
For every chunking of the same $(K,N,AD,M)$, the streaming sealer emits
the same $(C,T)$ as Algorithm~4.  For every chunking of the same
$(K,N,AD,C)$, the streaming opener accepts exactly when Algorithm~5
accepts and returns the same plaintext.
\end{proposition}

\begin{proof}[Proof sketch]
The subkeys, initial MAC state, and length counters are identical.  The
keystream state preserves the exact 64-byte block and offset across
calls; by induction on the chunk sequence, the concatenated XOR bytes
equal the one-shot ChaCha20 keystream with counters $0,1,\ldots$.  The
running MAC absorbs exactly $N\concat AD\concat C$ and then the two
length suffixes.  The opener performs the same operations, with the
additional invariant that tentative plaintext is zeroized on failure.
\end{proof}

\subsection{Failure rules}

\begin{table}[t]
\centering
\small
\begin{tabular}{@{}lll@{}}
\toprule
Condition & Sealer & Opener \\
\midrule
length mismatch & erase chunk, poison, $\bot$ & erase chunk, poison, $\bot$ \\
counter exhausted & erase chunk, poison, $\bot$ & erase chunk, poison, $\bot$ \\
length overflow & erase chunk, poison, $\bot$ & erase chunk, poison, $\bot$ \\
finish after poison & erase tag, $\bot$ & erase plaintext, $\bot$ \\
finish length mismatch & --- & erase plaintext, $\bot$ \\
tag mismatch & --- & erase plaintext, $\bot$ \\
success & write tag, $\top$ & validate, $\top$ \\
\bottomrule
\end{tabular}
\caption{Streaming AEAD failure rules.  $\top$ denotes \texttt{Some(())}
and $\bot$ denotes \texttt{None}.}
\label{tab:aeadfail}
\end{table}

A poisoned stream never recovers.  The opener's plaintext is tentative
until \texttt{finish} returns $\top$; callers MUST NOT parse, display,
execute, or transmit it before that point.

%======================================================================
\section{Signature specification}
\label{sec:ds}

\subsection{Parameters and sizes}

The construction follows the SPHINCS+/SLH-DSA family~\cite{sphincs,fips205};
XMSS is standardised in RFC~8391~\cite{rfc8391}.  HypG is not
wire-compatible with either standard.

\begin{table}[t]
\centering
\small
\begin{tabular}{@{}lr@{}}
\toprule
Symbol & Value \\
\midrule
$n$ & 32 bytes (BLAKE3-256) \\
$h$ & 68 \\
$d$ & 17 layers \\
$h'$ & 4 (16 leaves per XMSS tree) \\
$k_{\mathrm{FORS}}$ & 35 \\
$a$ & 9 (512 FORS leaves per tree) \\
$w$ & 16 \\
$\ell$ & 67 (64 message + 3 checksum digits) \\
$pk_{\mathrm{sig}}$ & 64 bytes \\
$sk_{\mathrm{sig}}$ & 128 bytes \\
$\mathit{sig}$ & 49,992 bytes \\
\bottomrule
\end{tabular}
\caption{Signature parameters.}
\label{tab:dsparams}
\end{table}

The sizes are:
\[
\begin{aligned}
 \mathrm{FORS\ unit} &= 32+9\cdot32=320,\\
 \mathrm{FORS\ total} &= 35\cdot320=11{,}200,\\
 \mathrm{XMSS\ layer} &= 8+67\cdot32+4\cdot32=2{,}280,\\
 \mathrm{hypertree} &= 17\cdot2{,}280=38{,}760,\\
 \mathit{sig} &= 32+11{,}200+38{,}760=49{,}992.
\end{aligned}
\]

\subsection{Address structure}

\begin{definition}[ADRS]
An address is 32 bytes:
\[
 \mathrm{layer}(4)\concat\mathrm{tree}(12)\concat
 \mathrm{type}(4)\concat f_{20}(4)\concat f_{24}(4)\concat f_{28}(4),
\]
all integers little-endian.  The 12-byte tree field holds a 96-bit
integer; valid hypertree addresses use at most 68 bits and zero the high
28 bits.
\end{definition}

\begin{table}[t]
\centering
\small
\begin{tabular}{@{}cllll@{}}
\toprule
Type & Name & $f_{20}$ & $f_{24}$ & $f_{28}$ \\
\midrule
0 & \texttt{T\_FORS} & FORS tree & FORS leaf & zero \\
1 & \texttt{T\_WOTS\_HASH} & keypair & chain & absolute step \\
2 & \texttt{T\_WOTS\_PK} & keypair & L-tree height & pair index \\
3 & \texttt{T\_TREE} & FORS tree or zero & node height & node index \\
4 & \texttt{T\_ROOT} & defined, unused & defined, unused & defined, unused \\
\bottomrule
\end{tabular}
\caption{Address types and field assignments.}
\label{tab:adrs}
\end{table}

\subsection{Tweakable hash functions}

All three functions use keyed BLAKE3-256 with the public seed
$\mathit{pubseed}$ as the 32-byte key:
\[
\begin{aligned}
 \PRF(\mathit{pubseed},x,A)
   &=\mathrm{KeyedBLAKE3}(\mathit{pubseed},
     01\concat x\concat A),\\
 F(\mathit{pubseed},A,x)
   &=\mathrm{KeyedBLAKE3}(\mathit{pubseed},
     02\concat A\concat x),\\
 \Hash(\mathit{pubseed},A,L,R)
   &=\mathrm{KeyedBLAKE3}(\mathit{pubseed},
     A\concat L\concat R).
\end{aligned}
\]
The first two inputs are 65 bytes; the third is 96 bytes.  The tags
$01$ and $02$ are the DS 2.0.1 revision.  $\Hash$ is untagged but its
arity and field order are distinct.

\subsection{Key generation}

\begin{steps}{Algorithm 6: $\mathrm{DS.KeyGen}(\mathit{esk})$}
\item $\mathit{pubseed}\gets\mathrm{CA\text{-}PRF}(\mathit{esk},
  \texttt{"HYP-G/v07/ds-pubseed"},32)$.
\item $\mathit{prf}\gets\mathrm{CA\text{-}PRF}(\mathit{esk},
  \texttt{"HYP-G/v07/ds-prf"},32)$.
\item $\mathit{root}\gets\mathrm{HT.Root}(\mathit{pubseed},
  \mathit{esk})$, the top XMSS layer at layer 16, tree 0.
\item $pk\gets\mathit{root}\concat\mathit{pubseed}$.
\item $sk\gets\mathit{esk}\concat\mathit{prf}\concat
  \mathit{pubseed}\concat\mathit{root}$.
\item Return $(pk,sk)$.
\end{steps}

\subsection{Randomizer, digest, and split}

\begin{steps}{Algorithm 7: $\mathrm{PreSign}(sk,\mathit{msg})$}
\item Parse $(\mathit{esk},\mathit{prf},\mathit{pubseed},
  \mathit{root})$ from $sk$.
\item $\mathit{stream}\gets\mathrm{CA\text{-}PRF}(\mathit{prf},
  \texttt{"HYP-G/v07/ds-addrnd"},64)$.
\item $R\gets\Hash(\mathit{stream}\concat\Hash(\mathit{msg}))$.
\item Compute, for $c\in\{0,1\}$,
  \[
   md_c=\Hash(R\concat\mathit{pubseed}\concat\mathit{root}
   \concat\Hash(\mathit{msg})\concat\mathrm{LE}_{32}(c)
   \concat\texttt{"HYP-G/v07/ds-msghash"}).
  \]
\item Split the first 387 bits of $md=md_0\concat md_1$:
  \[
   \mathit{fors}_t=\sum_{b=0}^{8}2^b\cdot
   \mathrm{bit}(md,9t+b),\quad t=0,\ldots,34,
  \]
  \[
   \mathit{tree}=\sum_{b=0}^{67}2^b\cdot
   \mathrm{bit}(md,315+b),\qquad
   \mathit{leaf}=\sum_{b=0}^{3}2^b\cdot
   \mathrm{bit}(md,383+b).
  \]
\item Return $(R,\mathit{fors},\mathit{tree},\mathit{leaf},
  \mathit{esk},\mathit{pubseed})$.
\end{steps}

\subsection{FORS}

\begin{definition}[FORS leaf and node]
For tree $t$ and leaf $l$:
\[
 A=\mathrm{ADRS}(0,0,0,t,l,0),\qquad
 \mathit{leaf}(t,l)=\PRF(\mathit{pubseed},\mathit{esk},A).
\]
For a node at height $h$ and pair index $p$:
\[
 A=\mathrm{ADRS}(0,0,3,t,h,p),\qquad
 \mathit{node}=\Hash(\mathit{pubseed},A,L,R).
\]
\end{definition}

The FORS public key is
\[
 \mathrm{FORS.PK}=\mathrm{KeyedBLAKE3}(\mathit{pubseed},
 A_{\mathrm{root}}\concat\mathit{root}_0\concat\cdots
 \concat\mathit{root}_{34}),
\]
where $A_{\mathrm{root}}=\mathrm{ADRS}(0,0,0,0xFFFFFFFF,0,0)$ and
$\mathit{root}_t$ is the root of tree $t$.  The input is 1152 bytes.

The signature for tree $t$ is the selected leaf plus its nine
authentication nodes.  Verification recomputes each root from the leaf
and path, then compresses the 35 roots.  An out-of-range selector
returns an all-zero unit; a real root cannot equal that output except
with negligible probability.

\begin{lemma}[FORS reconstruction]
\label{lem:fors}
Given a valid FORS signature for indices $\mathit{fors}$ and public
seed $\mathit{pubseed}$, $\mathrm{FORS.Verify}$ returns the same
$\mathrm{FORS.PK}$ as $\mathrm{FORS.PK}$ generation.
\end{lemma}

\begin{proof}
At height $h$, the verifier uses the current node index
$\lfloor\mathit{fors}_t/2^h\rfloor$, its sibling, and the address
$(t,h,\lfloor\mathit{fors}_t/2^{h+1}\rfloor)$.  The signer stored
exactly that sibling and the same address.  Induction on $h$ therefore
reconstructs each tree root.  Compression is deterministic.
\end{proof}

\subsection{WOTS+}

\begin{definition}[Digits]
For a 32-byte message:
\[
 d_i=(msg_{\lfloor i/2\rfloor}\gg4(i\bmod2))\land0x0F,
 \quad i=0,\ldots,63.
\]
The checksum is
\[
 c=\sum_{i=0}^{63}(15-d_i),
\]
encoded as three base-16 digits most-significant first:
$d_{64}=(c\gg8)\land0x0F$,
$d_{65}=(c\gg4)\land0x0F$,
$d_{66}=c\land0x0F$.
\end{definition}

\begin{definition}[Chain]
For a base address with layer, tree, and keypair index, and chain
index $i$:
\[
 \mathit{sk}_i=\PRF(\mathit{pubseed},\mathit{esk},
 \mathrm{ADRS}(\mathit{layer},\mathit{tree},1,i,0,0)).
\]
For a segment $[start,start+steps)$:
\[
 \mathrm{Chain}(x;start,steps,A)=
 F^{\,start+steps-1}_{A,\ldots}\circ\cdots\circ
 F^{\,start}_{A}(x),
\]
where each $F$ call uses address $A$ with $f_{24}=i$ and
$f_{28}=start+s$.  The implementation executes all 15 steps and
selects branchlessly.
\end{definition}

\begin{lemma}[WOTS chain compatibility]
\label{lem:wots-chain}
For $0\le d\le15$,
\[
 \mathrm{Chain}(\mathrm{Chain}(x;0,d,A);d,15-d,A)
 =\mathrm{Chain}(x;0,15,A).
\]
\end{lemma}

\begin{proof}
Both sides apply the same 15 chain steps with absolute addresses
$0,\ldots,14$; composition of functions is associative and the
intermediate segment boundary at $d$ does not change the addresses or
inputs.
\end{proof}

The L-tree compresses 67 chain ends.  At height $h$ and pair index $p$:
\[
 A=\mathrm{ADRS}(\mathit{layer},\mathit{tree},2,
 \mathit{keypair},h,p),\qquad
 \mathit{next}_p=\Hash(\mathit{pubseed},A,L,R).
\]
If the width is odd, the final unpaired node is promoted unchanged.
The widths are $67,34,17,9,5,3,2,1$.

\begin{lemma}[L-tree determinism]
\label{lem:ltree}
The L-tree output is a deterministic function of the 67 chain ends and
the base address.
\end{lemma}

\begin{proof}
Each level is a deterministic pairwise map plus a deterministic
promotion rule; induction on the width gives the claim.
\end{proof}

\subsection{XMSS}

An XMSS layer has 16 leaves.  For keypair $i$:
\[
 \begin{aligned}
 \mathit{base}&=\mathrm{ADRS}(\mathit{layer},\mathit{tree},0,i,0,0),\\
 \mathit{sk}&=\mathrm{WOTS.GenSK}(\mathit{pubseed},\mathit{esk},
   \mathit{base}),\\
 \mathit{pk}&=\mathrm{WOTS.PKFromSK}(\mathit{pubseed},\mathit{sk},
   \mathit{base}),\\
 \mathit{leaf}&=\mathrm{WOTS.LTree}(\mathit{pubseed},\mathit{pk},
   \mathit{base}).
 \end{aligned}
\]
The Merkle root at height $h$ and node index $p$ uses
\[
 A=\mathrm{ADRS}(\mathit{layer},\mathit{tree},3,0,h,p),
 \qquad
 \mathit{node}=\Hash(\mathit{pubseed},A,L,R).
\]
The XMSS signature is a WOTS+ signature plus four authentication
nodes.  Verification recomputes the leaf and climbs with the same
parity rule: if the current index is even, the computed node is the left
child; otherwise it is the right child.

\begin{lemma}[XMSS reconstruction]
\label{lem:xmss}
For a valid XMSS signature, verification returns the same root as
XMSS public-key generation.
\end{lemma}

\begin{proof}
The WOTS+ signature reconstructs the selected leaf by
Lemma~\ref{lem:wots-chain}, and the authentication path supplies the sibling at every
height.  Induction on height, with the address and parity rules above,
reconstructs the Merkle root.
\end{proof}

\subsection{Hypertree}

For layer $i=0,\ldots,16$:
\[
 \mathit{tree}_i=\mathit{tree\_idx}\gg4(i+1),\qquad
 \mathit{leaf}_0=\mathit{leaf},\qquad
 \mathit{leaf}_i=(\mathit{tree\_idx}\gg4i)\land15
 \quad(i\ge1).
\]
Layer 0 signs the FORS public key; each higher layer signs the root
returned by the layer below.  `HT.Root` is the XMSS root of layer 16,
tree 0.

\begin{lemma}[Hypertree consistency]
\label{lem:ht}
For an honestly generated key and signature, hypertree verification
returns `HT.Root`.
\end{lemma}

\begin{proof}
Induction on the layer index.  The base message is the FORS public key.
At layer $i$, XMSS reconstruction gives the next message; after layer
16 the message is the top root, which equals the public-key root by
construction.
\end{proof}

\subsection{Signing and verification}

\begin{steps}{Algorithm 8: $\mathrm{DS.Sign}(sk,msg)$}
\item $(R,\mathit{fors},\mathit{tree},\mathit{leaf},esk,pubseed)
  \gets\mathrm{PreSign}(sk,msg)$.
\item Compute the 35 FORS units by signing each $(\mathit{fors}_t)$.
\item Compute the FORS public key from the units (or from the secret).
\item $m_0\gets\mathrm{FORS.PK}$.
\item For $i=0,\ldots,16$:
  \begin{enumerate}[label=\alph*)]
  \item $s_i\gets\mathrm{XMSS.Sign}(\mathit{pubseed},esk,i,
    \mathit{tree}_i,\mathit{leaf}_i,m_i)$.
  \item $m_{i+1}\gets\mathrm{XMSS.Verify}(\mathit{pubseed},i,
    \mathit{tree}_i,\mathit{leaf}_i,m_i,s_i)$.
  \end{enumerate}
\item Serialize $R$, the 35 FORS units, and the 17 layer signatures,
  inserting eight zero bytes before each layer's WOTS+ signature.
\item Return the 49,992-byte signature.
\end{steps}

\begin{steps}{Algorithm 9: $\mathrm{DS.Verify}(pk,msg,sig)$}
\item If $|sig|\ne49{,}992$, return \texttt{false}.
\item Parse $R$, $\mathit{root}$, $\mathit{pubseed}$.
\item Recompute the digest and split into FORS indices, tree index, and
  leaf.
\item Parse the 35 FORS units and compute the FORS public key.
\item For each layer, require the eight reserved bytes to be zero.
\item Climb all 17 XMSS layers in order.
\item Return $\mathrm{ct\_eq}(\mathit{root}',\mathit{root})=1$.
\end{steps}

\subsection{Encoding}

\[
 \mathit{sig}=R(32)\concat
 \mathrm{FORS}_{35}\concat\mathrm{layer}_0\concat\cdots
 \concat\mathrm{layer}_{16}.
\]
Each FORS unit is $32+9\cdot32=320$ bytes.  Each layer is
\[
 \mathrm{zero}(8)\concat\mathrm{WOTS}(67\cdot32)
 \concat\mathrm{auth}(4\cdot32)=2280
\]
bytes.  The total is $32+35\cdot320+17\cdot2280=49{,}992$.

\subsection{DS 2.0.1 revision}

Relative to DS 2.0.0, revision 2.0.1 changes exactly two wire-visible
items:

\begin{enumerate}[leftmargin=*]
\item $\PRF$ prefixes its input with $01$; $F$ prefixes its input with
  $02$.
\item The WOTS+ L-tree promotes an unpaired final node unchanged
  instead of hashing it with itself.
\end{enumerate}

The first change removes the 64-byte structural collision between
$\PRF$ and $F$ inputs.  The second follows the FIPS 205 odd-node rule
and changes every WOTS+ L-tree root, every XMSS root, and every
signature.  The frozen file is \texttt{kats-ds/v2.0.1.txt}.  The
v2.0.0 DS vector file remains historical.

%======================================================================
\section{Domain registry}
\label{sec:domains}

\begin{longtable}{@{}lll@{}}
\caption{Domain registry.  Every active string is listed; the
\texttt{-init}/\texttt{-out} suffixes are CA-PRF mechanism, not
separate rows.}\label{tab:domains}\\
\toprule
Domain & Construction & Purpose \\
\midrule
\endhead
\texttt{HYP-G/v07/seed-a} & CA-PRF(seed, d, 32) & KEM $\mathit{seed}_A$ \\
\texttt{HYP-G/v07/sec-s} & CA-PRF(seed, d, 512) & KEM secret $s$ \\
\texttt{HYP-G/v07/sec-e} & CA-PRF(seed, d, 512) & KEM error $e$ \\
\texttt{HYP-G/v07/poly-a} & CA-PRF(seed$_A$, d, 6144) & matrix $A$ \\
\texttt{HYP-G/v07/rand-r} & CA-PRF(rseed, d, 512) & FO $r$ \\
\texttt{HYP-G/v07/rand-e1} & CA-PRF(rseed, d, 512) & FO $e_1$ \\
\texttt{HYP-G/v07/rand-e2} & CA-PRF(rseed, d, 128) & FO $e_2$ \\
\texttt{HYP-G/v07/fake} & CA-PRF(seed, d, 32) & rejection key \\
\texttt{HYP-G/v07/r-seed} & BLAKE3(m$\concat$pk$\concat$d) & FO seed \\
\texttt{HYP-G/v07/ss} & XOF(d$\concat$m$\concat$ct$\concat$H(pk),32) & accepted secret \\
\texttt{HYP-G/v07/reject} & XOF(d$\concat$fake$\concat$ct,32) & rejected secret \\
\texttt{HYP-G/v07/symm-enc} & BLAKE3 derive\_key & AEAD enc subkey \\
\texttt{HYP-G/v07/symm-mac} & BLAKE3 derive\_key & AEAD MAC subkey \\
\texttt{HYP-G/v07/ds-addrnd} & CA-PRF(prf, d, 64) & randomizer stream \\
\texttt{HYP-G/v07/ds-msghash} & digest suffix & digest separation \\
\texttt{HYP-G/v07/ds-pubseed} & CA-PRF(esk, d, 32) & public hash key \\
\texttt{HYP-G/v07/ds-prf} & CA-PRF(esk, d, 32) & randomizer seed \\
\texttt{HYP-G/v07/kat} & hash KAT only & never protocol use \\
\bottomrule
\end{longtable}

The CA-PRF mechanism appends \texttt{"-init"} and \texttt{"-out"} to a
domain; those suffixes are not separate registry rows.  The retired
draft strings \texttt{ds-fors}, \texttt{ds-wots}, and \texttt{ds-tree}
are not used.

%======================================================================
\section{API contracts}
\label{sec:api}

\subsection{KEM}

\begin{verbatim}
keygen(seed: &[u8; 32]) -> (PublicKey, SecretKey)
encaps(pk: &PublicKey, coins: &[u8; 32])
    -> (Ciphertext, SharedSecret)
decaps(sk: &SecretKey, ct: &Ciphertext) -> SharedSecret
\end{verbatim}
\texttt{decaps} never fails; an invalid ciphertext returns the
pseudorandom rejection secret.  \texttt{SecretKey} and
\texttt{SharedSecret} are not \texttt{Copy}; their \texttt{Debug} is
redacted; equality is constant-time; \texttt{Drop} best-effort
zeroizes.

\subsection{Hash and AEAD}

\begin{verbatim}
hash32, hash64, keyed32, derive32, expand, kdf, Hasher
seal(key, nonce, ad, pt, ct, tag) -> Option<()>
open(key, nonce, ad, ct, tag, pt) -> Option<()>
StreamSealer::new/update/finish
StreamOpener::new/update/finish
\end{verbatim}
One-shot \texttt{open} leaves the plaintext buffer untouched on
authentication failure.  Streaming opening is tentative until
\texttt{finish}.

\subsection{Signatures}

\begin{verbatim}
ds_keygen(esk: &[u8; 32]) -> (SigPk, SigSk)
ds_sign(sk: &SigSk, msg: &[u8], sig: &mut [u8; 49992])
ds_verify(pk: &SigPk, msg: &[u8], sig: &[u8]) -> bool
mt_sign(sk: &SigSk, msg: &[u8], sig: &mut [u8; SIG_LEN])
\end{verbatim}
\texttt{ds\_sign} is deterministic in \texttt{(sk,msg)}.
\texttt{ds\_verify} returns \texttt{false} on any length or encoding
error.  \texttt{mt\_sign} is byte-identical to \texttt{ds\_sign} and
never panics on a lost worker; the lost chunk is recomputed serially.

\subsection{Secret-lifetime contract}

\texttt{SecretKey}, \texttt{SharedSecret}, \texttt{SigSk},
\texttt{PreSign}, \texttt{Key}, \texttt{State}, \texttt{Rule}, and
\texttt{Hasher} redact \texttt{Debug} and best-effort zeroize on drop.
This is not volatile erasure.  Accessors such as
\texttt{seed\_copy}, \texttt{fake\_copy}, and \texttt{public\_copy}
return caller-owned copies; the caller must erase them.

%======================================================================
\section{Known-answer vectors}
\label{sec:kat}

\begin{table}[t]
\centering
\small
\begin{tabular}{@{}lrr@{}}
\toprule
File & Bytes & SHA-256 \\
\midrule
\texttt{kats/v2.0.0.txt} & 28,514 &
  \texttt{3b2dde3d...71d4b} \\
\texttt{kats-symm/v2.0.0.txt} & 3,662 &
  \texttt{91b7942f...4208} \\
\texttt{kats-hash/v2.0.0.txt} & 923 &
  \texttt{5fa77803...d9f1} \\
\texttt{kats-ds/v2.0.1.txt} & 201,015 &
  \texttt{e88c1bdf...c813} \\
\midrule
Total & 234,114 & \\
\bottomrule
\end{tabular}
\caption{Frozen files.  The complete hashes are in
\texttt{docs/parameters.md}.}
\label{tab:kat}
\end{table}

The complete SHA-256 values are:

\begin{verbatim}
kats/v2.0.0.txt
3b2dde3d1138c11b54a6880f50ca7b4390071ae53aa94601075f39b2af971d4b
kats-symm/v2.0.0.txt
91b7942ffb4a3c5da8b8b589268af7b7d1594796b428efdc7e26c567c5f04208
kats-hash/v2.0.0.txt
5fa778035ce2e8aa9c89345458ba4da67df4bdbe71a456cbb58026cf9341d9f1
kats-ds/v2.0.1.txt
e88c1bdf54e9390736cd03c1a165616d744a42c48d662880b344c7de3546c813
\end{verbatim}

The KEM file contains three deterministic tag groups
($00$, $42$, $FF$) and a BLAKE3 \texttt{"abc"} anchor.  The Symm file
contains five AEAD cases (empty, one-byte, 127/128-byte, and 1000-byte
messages with varied AD).  The Hash file contains empty, \texttt{"abc"},
and 1025-byte pattern cases plus keyed and derive-key anchors.  The DS
file contains two cases: an empty message and a 55-byte message, each
with a 49,992-byte signature.

The x86-64 and \texttt{qemu-aarch64} musl-static KAT outputs are
byte-identical across all 234,114 bytes.  This checks endianness and
integer width on an emulated ARM64 target; it is not native ARM64
validation.

%======================================================================
\section{Validation matrix}
\label{sec:validation}

\begin{table}[t]
\centering
\small
\begin{tabular}{@{}lll@{}}
\toprule
Gate & Result & Evidence class \\
\midrule
Debug tests & 146/146 & E1 \\
Release tests & 146/146 & E1 \\
MSRV 1.70 tests & 146/146 & E1 \\
Red-team S1--S14 & 14/14 & E2 \\
Failure sweep (release) & 20,000/20,000 & E2 \\
Full-bit flip sweep (one ct) & 11,520/11,520 & E2 \\
Lattice lab L1--L5 & 5/5 & E2 \\
Cross-architecture KAT & 234,114 bytes identical & E3 \\
Miri KEM & 9/9 & E3 \\
Miri AEAD & 6/6 & E3 \\
Miri DS bounds & 3/3 & E3 \\
\bottomrule
\end{tabular}
\caption{Validation record.  E4 (independent audit) is absent.}
\label{tab:validation}
\end{table}

The debug, release, and MSRV suites each contain 108 unit tests, 25 CLI
negative integration tests, and 13 doc-tests.  Full DS signing is not
covered by Miri.  The failure sweep and bit-flip sweep are sanity
checks, not a formal failure bound.  The lattice lab is not a BKZ cost
estimate.

\begin{verbatim}
cargo test --workspace --locked
cargo test --workspace --locked --release
cargo +1.70 test --workspace --locked
cargo run -q -p hypg-redteam
cargo run -q -p hypg-lattice
cargo test -p hypg-kat --lib --locked
\end{verbatim}

%======================================================================
\section{Conformance and change control}
\label{sec:conformance}

An implementation conforms to this document for a given revision only
if it reproduces the active KAT files byte-for-byte and satisfies the
specified failure and encoding rules.  A change is wire-visible if it
alters any of:

\begin{itemize}[leftmargin=*]
\item a ring parameter, modulus, distribution, packing, compression, or
  message encoding;
\item a domain string, tag byte, hash input order, or ADRS field;
\item a DS tree shape, digest split, checksum rule, or reserved byte;
\item an AEAD key-derivation context, MAC input, or length limit;
\item a frozen KAT file or its interpretation.
\end{itemize}

Any such change requires a new version identifier, regenerated KATs, an
updated registry, a cross-architecture equality run, and a new
specification revision.  A test, benchmark, or zeroization change that
does not alter bytes is not wire-visible, but must still preserve the
frozen output.

%======================================================================
\section{Security status and non-claims}
\label{sec:security}

This specification does not prove any security property.  The intended
analysis depends on Module-LWE hardness, BLAKE3, ChaCha20, the exact
FO-style transform, random-oracle behaviour, a negligible decapsulation
failure rate, nonce uniqueness, and a bounded signing volume.
Published FO analyses~\cite{fo99,hhk17} do not automatically cover the
chosen accepted-secret and rejection derivations.  None of
these is established by the implementation or by this document.

The KEM has no exact published proof for its shared-secret and
rejection derivations; the lattice-estimator has not been run; the
signature composition has no exact UF-CMA/SUF-CMA proof or
signing-volume bound; the AEAD has no reduction for the exact
composition; and no external audit, formal verification, constant-time
machine-code proof, or power/EM/fault campaign exists.  HypG is a
research prototype and MUST NOT be deployed to protect real assets.

%======================================================================
\appendix
\section{Constant inventory}
\label{app:constants}

\begin{longtable}{@{}ll@{}}
\caption{Constant inventory.}\\
\toprule
Constant & Value \\
\midrule
\endhead
\texttt{STATE\_WORDS} & 32 \\
\texttt{STATE\_BYTES} & 256 \\
\texttt{RULE\_WORDS} & 8 \\
\texttt{ROUNDS} & 32 \\
\texttt{SEED\_LEN} & 32 \\
\texttt{FAKE\_LEN} & 32 \\
\texttt{PK\_LEN} & 1568 \\
\texttt{CT\_LEN} & 1440 \\
\texttt{SK\_LEN} & 1632 \\
\texttt{SS\_LEN} & 32 \\
\texttt{POLY\_N} & 256 \\
\texttt{MODULE\_K} & 4 \\
\texttt{MOD\_Q} & 4096 \\
\texttt{C1\_Q} & 1024 \\
\texttt{C2\_Q} & 32 \\
\texttt{CBD\_ETA} & 2 \\
\texttt{POLY\_PACKED\_LEN} & 384 \\
\texttt{C1\_PACKED\_LEN} & 320 \\
\texttt{C2\_PACKED\_LEN} & 160 \\
\texttt{SIG\_PK\_LEN} & 64 \\
\texttt{SIG\_SK\_LEN} & 128 \\
\texttt{SIG\_LEN} & 49,992 \\
\bottomrule
\end{longtable}

\section{Byte-offset quick reference}
\label{app:offsets}

\begin{center}
\small
\begin{tabular}{@{}rrl@{}}
\toprule
Object & Offset & Field \\
\midrule
$sk$ & 0 & master seed (32) \\
$sk$ & 32 & fake key (32) \\
$sk$ & 64 & public key (1568) \\
$sig$ & 0 & $R$ (32) \\
$sig$ & 32 & FORS units (11,200) \\
$sig$ & 11,232 & layer 0 (2,280) \\
$sig$ & 47,712 & layer 16 (2,280) \\
\bottomrule
\end{tabular}
\end{center}
\noindent\textbf{Selected byte offsets.}  Each layer begins with eight
zero bytes, followed by 67 WOTS+ values and four authentication nodes.

\section{CLI and error codes}
\label{app:cli}

The operator CLI is not a production key-management tool.  Exit code
$0$ means success; code $2$ means usage or input error; code $3$ means
authentication failure (bad AEAD tag or bad signature).  KEM
decapsulation of an invalid ciphertext exits $0$ because the library
returns a pseudorandom secret rather than an error.

\section{Comparison with FIPS 203 and FIPS 205}
\label{app:fips}

HypG resembles ML-KEM-1024~\cite{fips203} in module rank and ring
degree but uses $q=4096$, 10-bit $c_1$ compression, a different
FO-style secret derivation, and BLAKE3.  It resembles SLH-DSA-256f
(FIPS~205)~\cite{fips205} in its tree dimensions but uses BLAKE3, a
custom 32-byte address, explicit $\PRF$/$F$ tags, deterministic
randomizer derivation, and eight reserved bytes per layer.  It is not
ML-KEM, not SLH-DSA, and not FIPS-conformant.  Matching dimensions is
not a security certification.  ML-DSA (FIPS~204)~\cite{fips204} is not
used by HypG.

\begin{thebibliography}{99}

\bibitem{rfc2119}
S.~Bradner.
\newblock Key words for use in RFCs to Indicate Requirement Levels.
\newblock RFC 2119, March 1997.

\bibitem{fips203}
National Institute of Standards and Technology.
\newblock \emph{Module-Lattice-Based Key-Encapsulation Mechanism
Standard}. FIPS 203, August 2024.

\bibitem{fips204}
National Institute of Standards and Technology.
\newblock \emph{Module-Lattice-Based Digital Signature Standard}.
FIPS 204, August 2024.

\bibitem{fips205}
National Institute of Standards and Technology.
\newblock \emph{Stateless Hash-Based Digital Signature Standard}.
FIPS 205, August 2024.

\bibitem{blake3}
J.~O'Connor, J.-P. Aumasson, S.~Neves, and Z.~Wilcox-O'Hearn.
\newblock BLAKE3: one function, fast everywhere.
\newblock Technical report, 2020.

\bibitem{rfc8439}
Y.~Nir and A.~Langley.
\newblock ChaCha20 and Poly1305 for IETF Protocols.
\newblock RFC 8439, June 2018.

\bibitem{rfc8391}
A.~H\"ulsing, D.~Butin, S.-L. Gazdag, J.~Rijneveld, and A.~Mohaisen.
\newblock XMSS: eXtended Merkle Signature Scheme.
\newblock RFC 8391, May 2018.

\bibitem{sphincs}
D.~J. Bernstein, A.~H\"ulsing, S.~K\"olbl, R.~Niederhagen,
J.~Rijneveld, and P.~Schwabe.
\newblock The SPHINCS+ signature framework.
\newblock In \emph{ACM CCS 2019}, pages 2129--2146, 2019.

\bibitem{fo99}
E.~Fujisaki and T.~Okamoto.
\newblock Secure integration of asymmetric and symmetric encryption
schemes.
\newblock In \emph{CRYPTO 1999}, LNCS 1666, pages 537--554, 1999.

\bibitem{hhk17}
D.~Hofheinz, K.~H\"ovelmanns, and E.~Kiltz.
\newblock A modular analysis of the Fujisaki--Okamoto transformation.
\newblock In \emph{TCC 2017}, LNCS 10677, pages 341--371, 2017.

\end{thebibliography}

\end{document}
